%!TEX root = ../main/main.tex

Responda las siguientes preguntas.
\begin{enumerate}
\item \text{[0.5 puntos]} Para el caso de cuatro participantes, demuestre que si el participante 3 construye una firma
$(s_1, s_2, s_3, s_4, c_1, f)$ de un mensaje $m$ utilizando el protocolo, entonces el proceso
de verificación concluye que la firma es válida.

\item \text{[0.5 puntos]} Un participante $i$ podría tratar de ocultar que generó dos firmas utilizando $f = T^{x_i}$ en una de las
firmas, y un elemento arbitrario $f' \in G$ con $f' \neq f$ en la otra firma. Considerando nuevamente el caso con cuatros
participantes, explique por qué el participante 1 no podría generar una firma válida si utiliza un elemento $f' \in G$
tal que $f' \neq T^{x_1}$.

\item \text{[0.5 puntos]} Demuestre que $|\langle T \rangle| = q$.

\item \text{[0.5 puntos]} Suponga que modificamos el protocolo de manera que el elemento $T$ usado en el protocolo es generado por
una tercera parte; por ejemplo, este elemento $T$ podría ser generado por el Servel en caso de una elección o por los
profesores del curso en el caso de esta tarea. Muestre que, en este caso, el protocolo puede ser quebrado. En particular,
muestre que esta tercera parte puede generar eficientemente un elemento $T \in G$ tal que $|\langle T \rangle| = q$, $T \neq g$
y, si $T$ se utiliza en la generación de firmas de anillo, entonces puede determinar eficientemente quién
generó cada firma.
\end{enumerate}

\medskip
